\documentclass[12pt,a4paper]{article}

% --------------------------------------------------
% Packages
% --------------------------------------------------
\usepackage{amsmath,amssymb,mathtools}
\usepackage{physics}
\usepackage{graphicx}
\usepackage{booktabs}
\usepackage{array}
\usepackage{siunitx}
\usepackage{xcolor}
\usepackage{hyperref}
\usepackage{enumitem}
\usepackage{fourier}
\usepackage{charter}


\usepackage[margin=2.5cm]{geometry}
\title{A Sample Physics Document}
\author{Baijayanta Bhattacharyya}
\date{\today}

\begin{document}
	
	\maketitle
	
	\begin{abstract}
		This is a sample LaTeX document containing mathematical equations,
		tables, matrices, figures, lists, and references. It can be used to
		test conversion from \LaTeX{} to HTML.
	\end{abstract}
	
	\tableofcontents
	
	% ==================================================
	\section{Introduction}
	% ==================================================
	
	Classical mechanics provides a powerful framework for describing
	physical systems. The motion of a system can be obtained from the
	Lagrangian
	
	\begin{equation}
		L = T - V,
		\label{eq:lagrangian}
	\end{equation}
	
	where $T$ is the kinetic energy and $V$ is the potential energy.
	
	The Euler--Lagrange equation is
	
	\begin{equation}
		\frac{d}{dt}
		\left(
		\frac{\partial L}{\partial \dot{q}_i}
		\right)
		-
		\frac{\partial L}{\partial q_i}
		= 0.
		\label{eq:euler-lagrange}
	\end{equation}
	
	Equation~\eqref{eq:euler-lagrange} determines the equations of motion
	of the system.
	
	% ==================================================
	\section{Mathematical Expressions}
	% ==================================================
	
	Consider a simple harmonic oscillator with Lagrangian
	
	\begin{equation}
		L =
		\frac{1}{2}m\dot{x}^{\,2}
		-
		\frac{1}{2}kx^2.
	\end{equation}
	
	Using the Euler--Lagrange equation,
	
	\begin{align}
		\frac{\partial L}{\partial x}
		&= -kx, \\
		\frac{\partial L}{\partial \dot{x}}
		&= m\dot{x}.
	\end{align}
	
	Therefore,
	
	\begin{equation}
		m\ddot{x} + kx = 0.
	\end{equation}
	
	Defining
	
	\begin{equation}
		\omega = \sqrt{\frac{k}{m}},
	\end{equation}
	
	we obtain
	
	\begin{equation}
		\boxed{
			\ddot{x} + \omega^2 x = 0
		}.
	\end{equation}
	
	The general solution is
	
	\begin{equation}
		x(t)
		=
		A\cos(\omega t)
		+
		B\sin(\omega t).
	\end{equation}
	
	% ==================================================
	\subsection{A More Complicated Equation}
	% ==================================================
	
	Consider the integral
	
	\begin{equation}
		I =
		\int_0^\infty
		\frac{\ln x}{1+x^3}\,dx.
	\end{equation}
	
	A useful identity is
	
	\begin{equation}
		\int_0^\infty
		\frac{x^{a-1}}{1+x^b}\,dx
		=
		\frac{\pi}{b}
		\csc\left(\frac{\pi a}{b}\right),
		\qquad
		0<a<b.
	\end{equation}
	
	Differentiating with respect to $a$ gives
	
	\begin{equation}
		\frac{d}{da}
		\int_0^\infty
		\frac{x^{a-1}}{1+x^b}\,dx
		=
		\int_0^\infty
		\frac{x^{a-1}\ln x}{1+x^b}\,dx.
	\end{equation}
	
	% ==================================================
	\section{Matrices}
	% ==================================================
	
	A general $2\times2$ matrix can be written as
	
	\begin{equation}
		A =
		\begin{pmatrix}
			a & b \\
			c & d
		\end{pmatrix}.
	\end{equation}
	
	Its determinant is
	
	\begin{equation}
		\det A = ad-bc.
	\end{equation}
	
	The inverse, provided $\det A\neq0$, is
	
	\begin{equation}
		A^{-1}
		=
		\frac{1}{ad-bc}
		\begin{pmatrix}
			d & -b \\
			-c & a
		\end{pmatrix}.
	\end{equation}
	
	For angular momentum, the commutation relation is
	
	\begin{equation}
		[J_i,J_j]
		=
		i\hbar\epsilon_{ijk}J_k.
	\end{equation}
	
	For spin-$\frac12$, the generators are
	
	\begin{equation}
		J_i = \frac{\hbar}{2}\sigma_i,
	\end{equation}
	
	where the Pauli matrices are
	
	\begin{align}
		\sigma_x &=
		\begin{pmatrix}
			0 & 1 \\
			1 & 0
		\end{pmatrix},
		&
		\sigma_y &=
		\begin{pmatrix}
			0 & -i \\
			i & 0
		\end{pmatrix},
		&
		\sigma_z &=
		\begin{pmatrix}
			1 & 0 \\
			0 & -1
		\end{pmatrix}.
	\end{align}
	
	% ==================================================
	\section{Tables}
	% ==================================================
	
	Table~\ref{tab:particles} shows some example particle properties.
	
	\begin{table}[ht]
		\centering
		\caption{Properties of selected particles.}
		\label{tab:particles}
		
		\begin{tabular}{lccc}
			\toprule
			Particle & Mass (\si{\kilo\gram}) &
			Charge ($e$) & Spin \\
			\midrule
			Electron & $9.11\times10^{-31}$ & $-1$ &
			$\frac12$ \\
			Proton   & $1.67\times10^{-27}$ & $+1$ &
			$\frac12$ \\
			Neutron  & $1.67\times10^{-27}$ & $0$ &
			$\frac12$ \\
			Photon   & $0$ & $0$ & $1$ \\
			\bottomrule
		\end{tabular}
	\end{table}
	
	% ==================================================
	\section{Lists}
	% ==================================================
	
	Important concepts in classical mechanics include:
	
	\begin{itemize}
		\item Configuration space
		\item Generalized coordinates
		\item Lagrangian mechanics
		\item Hamiltonian mechanics
		\item Canonical transformations
	\end{itemize}
	
	The usual procedure for solving a Lagrangian problem is:
	
	\begin{enumerate}
		\item Choose generalized coordinates.
		\item Write the kinetic energy.
		\item Write the potential energy.
		\item Construct the Lagrangian.
		\item Apply the Euler--Lagrange equations.
	\end{enumerate}
	
	% ==================================================
	\section{Figures}
	% ==================================================
	
	Figure~\ref{fig:example} is an example figure.
	
	\begin{figure}[ht]
		\centering
		\includegraphics[width=0.65\textwidth]{example-image-c}
		\caption{An example figure.}
		\label{fig:example}
	\end{figure}
	
	% ==================================================
	\section{Physics Example}
	% ==================================================
	
	Consider a simple pendulum of length $L$ and mass $m$.
	
	The position of the bob can be written as
	
	\begin{equation}
		x = L\sin\theta,
		\qquad
		y = -L\cos\theta.
	\end{equation}
	
	The velocity is
	
	\begin{equation}
		v^2 = L^2\dot{\theta}^{\,2}.
	\end{equation}
	
	Therefore, the kinetic and potential energies are
	
	\begin{align}
		T &= \frac12 mL^2\dot{\theta}^{\,2},\\
		V &= mgL(1-\cos\theta).
	\end{align}
	
	Hence,
	
	\begin{equation}
		L
		=
		\frac12mL^2\dot{\theta}^{\,2}
		-
		mgL(1-\cos\theta).
	\end{equation}
	
	Applying the Euler--Lagrange equation gives
	
	\begin{equation}
		mL^2\ddot{\theta}
		+
		mgL\sin\theta
		=0.
	\end{equation}
	
	Thus,
	
	\begin{equation}
		\boxed{
			\ddot{\theta}
			+
			\frac{g}{L}\sin\theta
			=0
		}.
		\label{eq:pendulum}
	\end{equation}
	
	For small oscillations,
	
	\begin{equation}
		\sin\theta \approx \theta,
	\end{equation}
	
	and therefore
	
	\begin{equation}
		\ddot{\theta}
		+
		\frac{g}{L}\theta
		=0.
	\end{equation}
	
	The period is
	
	\begin{equation}
		T = 2\pi\sqrt{\frac{L}{g}}.
	\end{equation}
	
	% ==================================================
	\section{References}
	% ==================================================
	
	For more information, see standard texts on classical mechanics
	such as~\cite{goldstein}.
	
	\begin{thebibliography}{9}
		
		\bibitem{goldstein}
		H. Goldstein, C. Poole, and J. Safko,
		\textit{Classical Mechanics},
		3rd ed.,
		Pearson, 2002.
		
		\bibitem{griffiths}
		D. J. Griffiths,
		\textit{Introduction to Quantum Mechanics},
		3rd ed.,
		Cambridge University Press, 2018.
		
	\end{thebibliography}
	
	% ==================================================
	\section{Conclusion}
	% ==================================================
	
	This document contains:
	
	\begin{itemize}
		\item Mathematical equations
		\item Aligned equations
		\item Matrices
		\item Tables
		\item Figures
		\item Lists
		\item Cross-references
		\item Citations
		\item Sections and subsections
	\end{itemize}
	
	These features make it useful for testing a \LaTeX{} to HTML conversion
	workflow.
	
\end{document}

